\documentclass{article}
\usepackage{amsmath,amssymb}

\title{Gifford-McMahon Cryocooler Design \& Build Notebook}
\author{Ira Parsons (notetaker) \& Chas Skidmore}
\date{Last Updated: \today}

\begin{document}

\maketitle

\section*{September 9, 2026}

\subsection*{Criteria:}

\begin{itemize}
\item Reach -200 C / 73 K or lower
\end{itemize}

\subsection*{Constraints:}

\begin{itemize}
 \item Run off wall outlet power
 \item Cold head outer diameter less than 3"
\end{itemize}

\subsection*{Material Selection}

Main body stainless steel.  For insulation, fiberglass wool, ceramic wool, or other?  Cold head heat exchanger copper or brass.  Brass pneumatic fittings.  For regenerator contents, metal beads vs. mesh, what metal has highest $c$ at planned operating temp (copper, brass, lead, stainless steel)?

At 75 K, copper has $c_p = 11.89 \text{ } \text{J}\cdot\text{K}^{-1}\cdot\text{mol}^{-1}$ (https://tsapps.nist.gov/srmext/certificates/archives/5.pdf).  At 75 K, 302 Stainless Steel has $c = 197.08013 \text{ J/(kg-K)}$ (https://trc.nist.gov/cryogenics/materials/304Stainless/304Stainless\_rev.htm).

Water cooling?

Main dimensions?  Regenerator length to cold chamber length ratio best around 6.5 (https://www.osti.gov/servlets/purl/14637).

\end{document}